\scanpage{11}

% loc: top continuation
munie d'un isom. $\alpha$ rendant commutatif

% loc: top diagram
\diagram{\begin{tikzcd}[ampersand replacement=\&,row sep=small]
\deleted{\unsure{\Phi\,\Pi\,\Phi(e)}}\&\\
\Phi_1(e)\arrow[r,shift left=.9ex,"\Pi_1", ""{name=A,below}]\arrow[r,bend right=16, "{\text{can}^{\unsure{\mathrm{on}|\circ}}}"',""{name=B,above}]\&C\arrow[Rightarrow,from=A,to=B,"\alpha"]\\[-.7em]
\deleted{\unsure{\gamma}}\arrow[u]\arrow[r,dash]\&\deleted{\unsure{e}}
\end{tikzcd}}

% loc: definition
\uline{Définition}\quad Si $(I,(L_i)_{i\in I})$ est une famille finie d'objets de $C$, $\Pi((L_i))$ est appelé le produit $\otimes$ de la famille dans $C$.

% loc: converse
Il y a une réciproque : Si on se donne \deleted{\unsure{une catégorie}} \deleted{\unsure{avec}}\inserted{une cat.} \inserted{avec} une opération $\otimes$, alors la donnée d'une structure $\unsure{\mathrm{AUC}|\mathrm{ANC}}$ sur celle-ci est \unsure{équivalent|équivalente} à la donnée d'un foncteur $\Pi:\Phi(e)\to C$, \deleted{et d'un isom. $\Pi(X\otimes Y)\simeq\Pi(X)\otimes\Pi(Y)$, en associant à tte structure d'associativité} qui sur $\Phi_1(C)$ soit ``l'identité'', et \deleted{\unsure{sur}} qui \deleted{\unsure{sur}} rend commutatif de plus le diagramme

% loc: middle triangle
\diagram{\begin{tikzcd}[ampersand replacement=\&]\Phi_2(C)\arrow[r,"\Pi_2"]\& C\\C\times C\arrow[u]\arrow[ur,"\otimes"']\&\end{tikzcd}}

% loc: ruled separator
\rule{2.3em}{1.4pt}

% loc: NB
\uline{NB}\quad Dans $\Phi(C)$ lui-m\textsuperscript{e}, si on a une famille finie d'objets de $\Phi(C)$, $(\xi_j)_{j\in J}$, $\xi_j=(I_j,(L_{ji})_{i\in I_j})$, alors $\Pi((\xi_j)_{j\in J})$ est

% loc: NB display
\[\left(\coprod_{j\in J}I_j,\ (L_{ij})_{\substack{j\in J\\i\in I_j}}\right)\quad\text{(avec abus de \uline{langage})}\]

% loc: foot paragraph
Revenons à la $\otimes$-\unsure{ct|cat} $C$ AUC. Soit $I$ un ens. \deleted{\unsure{(I)}} à \uline{deux} él\textsuperscript{nts}, $(L_i)_{i\in I}$ une famille d'objets de $C$ indexée par $I$. Il y a deux ordres totaux sur $I$, et si les deux él\textsuperscript{nts} de $I$ sont notés $a,b$, alors les deux façons de représenter le produit, suivant l'ordre total choisi, sont $L_a L_b$ et $L_b L_a$, on passe de l'un à l'autre par l'isom de symétrie. Cela offre \deleted{\unsure{deux}} du nouveau si $L_a=L_b=L$, car alors on trouve un \deleted{\unsure{isom} $s(L)$ \unsure{de L sur lui}} autom. d'ordre 2 de $L\otimes L$ \deleted{\unsure{qui}}

% archive sheet number: 10
